\section{Motivation: can the Artin--Schreier calculation be iterated?}
\label{sec:as_induction_motivation}

Part~(2) of \Cref{theorem:cyclic_cases_after_blowup} proves the order-$p$
case by an explicit Artin--Schreier calculation. Part~(3) treats
$G=\Z/p^r\Z$ by working directly with an Artin--Schreier--Witt equation.
It is natural to ask whether the second calculation can instead be obtained
by iterating the first.

Indeed, let
\[
G=\Z/p^r\Z,\qquad
H=p^{r-1}G\cong\Z/p\Z,\qquad
G/H\cong\Z/p^{r-1}\Z,
\]
and factor a $G$-torsor $\cP\to U$ as
\[
\cP\longrightarrow\cP/H\longrightarrow U.
\]
The tempting induction has three steps:
\begin{enumerate}
  \item apply the induction hypothesis to the $G/H$-torsor
        $\cP/H\to U$;
  \item regard $\cP\to\cP/H$ as an $H\cong\Z/p\Z$-torsor and apply the
        Artin--Schreier calculation after compactifying its base;
  \item combine the two unramifiedness statements to recover the result for
        the original $G$-torsor.
\end{enumerate}

The first two steps do produce natural statements, but they live over
different exceptional valuations. Passing from one to the other involves a
base change induced by the comparison between the symmetric power of a
torsor as a scheme and its symmetric power as a torsor. That base change can
be wildly ramified. The purpose of this appendix is to formulate the proposed
induction precisely, follow both inputs through the resulting valuation
tower, and explain why they do not by themselves imply the required
conclusion.


\section{Preliminaries for the attempted induction}
\label{sec:as_induction_preliminaries}

\subsection{The symmetric quotient tower}
\label{subsection:quotient_torsors_towers}

Let $G$ be a finite abelian group, let $H\subseteq G$ be a subgroup, and let
$\cP\to X$ be a $G$-torsor over an $S$-scheme $X$ that is Zariski-locally
quasi-projective over $S$. Recall the canonical factorization
\[
\cP\longrightarrow\cP/H\longrightarrow X,
\]
where $\cP/H=(G/H)\times^G\cP$ is a $G/H$-torsor over $X$ and
$\cP\to\cP/H$ is an $H$-torsor.

In addition to the kernel $\KG=\ker(G^d\to G)$ of the summation morphism,
put
\[
K_H:=\ker(H^d\to H),\qquad
K_{G/H}:=\ker((G/H)^d\to G/H).
\]
The snake lemma applied to the two summation maps gives an exact sequence
\begin{equation}\label{eq:summation_kernels_ses}
0\longrightarrow K_H\longrightarrow\KG\longrightarrow K_{G/H}
\longrightarrow0.
\end{equation}

\begin{construction}[The symmetric-power comparison morphism]
\label{construction:symmetric_power_comparison}
For the $G$-torsor $\cP\to X$ fixed above, put
\[
\Xi_G:=\KG\rtimes S_d.
\]
The inclusion $S_d\hookrightarrow\Xi_G$, given by
$\sigma\mapsto((1,\ldots,1),\sigma)$, induces a canonical comparison
morphism
\[
q_{\cP}:\cP^{(d)}
=S_d\backslash\cP^{\times_S d}
\longrightarrow
\Xi_G\backslash\cP^{\times_S d}
=\cP^{[d]}.
\]
Indeed, the quotient map from $\cP^{\times_S d}$ to its $\Xi_G$-quotient
is $S_d$-invariant and therefore factors uniquely through the
$S_d$-quotient.
\end{construction}

We use the following four quotient schemes:
\begin{align*}
  \cP^{[d]_G}&:=(\KG\rtimes S_d)\backslash\cP^{\times_S d},\\
  \cP^{[d]_H}&:=(K_H\rtimes S_d)\backslash\cP^{\times_S d},\\
  (\cP/H)^{(d)}&:=(H^d\rtimes S_d)\backslash\cP^{\times_S d},\\
  (\cP/H)^{[d]}&:=((H^d+\KG)\rtimes S_d)
      \backslash\cP^{\times_S d}.
\end{align*}
The first is the symmetric $G$-power of $\cP$. The second is the symmetric
$H$-power relative to $\cP\to\cP/H$, the third is the symmetric power of the
scheme $\cP/H$, and the fourth is the symmetric $G/H$-power of the
$G/H$-torsor $\cP/H\to X$. For the last two identifications, use
$(\cP/H)^{\times_S d}\cong H^d\backslash\cP^{\times_S d}$ and note that the
preimage of $K_{G/H}$ under $G^d\twoheadrightarrow(G/H)^d$ is
$H^d+\KG$.

\begin{prop}\label{prop:torsor_tower_diagram}
The four quotient schemes fit into a canonical commutative diagram
\[
\begin{tikzcd}[row sep=large, column sep=large]
  \cP^{[d]_H}\arrow[r,"H"]\arrow[d]
    &(\cP/H)^{(d)}\arrow[d,"q"]\\
  \cP^{[d]_G}\arrow[r,"H"']\arrow[dr,"G"',bend right=20]
    &(\cP/H)^{[d]}\arrow[d,"G/H"]\\
    &X^{(d)},
\end{tikzcd}
\]
in which every labelled arrow is a torsor under the indicated constant
group. The upper square is Cartesian:
\[
\cP^{[d]_H}\cong
\cP^{[d]_G}\times_{(\cP/H)^{[d]}}(\cP/H)^{(d)}.
\]
The vertical map $q=q_{\cP/H}$ is the comparison morphism of
\Cref{construction:symmetric_power_comparison}, applied to the
$G/H$-torsor $\cP/H\to X$.
\end{prop}

\begin{proof}
Apply \Cref{prop:symmetric_power_torsor} to the $G$-torsor $\cP\to X$, the
$H$-torsor $\cP\to\cP/H$, and the $G/H$-torsor $\cP/H\to X$. This gives the
three labelled torsors; the comparison map $q$ is supplied by
\Cref{construction:symmetric_power_comparison}, applied to
$\cP/H\to X$. The factorization
\[
\cP^{[d]_G}\longrightarrow(\cP/H)^{[d]}\longrightarrow X^{(d)}
\]
is the canonical quotient decomposition of the $G$-torsor
$\cP^{[d]_G}\to X^{(d)}$ by $H$, because both its intermediate quotient and
$(\cP/H)^{[d]}$ are realized as
\[
((H^d+\KG)\rtimes S_d)\backslash\cP^{\times_S d}.
\]

Let
\[
T:=\cP^{[d]_G}\times_{(\cP/H)^{[d]}}(\cP/H)^{(d)}.
\]
Base change makes $T$ an $H$-torsor over $(\cP/H)^{(d)}$. The quotient maps
give an $H$-equivariant morphism $\cP^{[d]_H}\to T$. A morphism of
$H$-torsors is an isomorphism, proving the Cartesian assertion.
\end{proof}

\begin{rem*}
The two vertical comparison maps in the diagram need not be torsors under a
constant group. For example, $K_H\rtimes S_d$ need not be normal in
$\KG\rtimes S_d$, because coordinate permutations act nontrivially on
$\KG/K_H$. Thus the map $q$ is not another torsor layer to which the cyclic
induction hypotheses automatically apply. Its ramification at the boundary
is precisely the extra datum that will matter below.
\end{rem*}


\subsection{The geometric boundary and its valuations}

Let $C$ be a smooth projective curve, let $U=C\setminus\mdls$, and let
$\cP\to U$ be a $G$-torsor. Fix $p_\infty\in\mdls$ and an integer $d\geq1$,
and assume that the ramification of $\cP$ at $p_\infty$ is bounded by $d$ in
the sense of \Cref{definition:local_ramification_boundness}.

The $G/H$-torsor $\cP/H\to U$ compactifies to a finite morphism of smooth
projective curves
\[
f:C'\longrightarrow C.
\]
Choose $p'_\infty\in f^{-1}(p_\infty)$. The induced morphism on symmetric
powers sends $d\cdot p'_\infty$ to $d\cdot p_\infty$:
\[
f^{(d)}:C'^{(d)}\longrightarrow C^{(d)}.
\]
Let
\[
\pi_C:X_C\longrightarrow C^{(d)},\qquad
\pi_{C'}:X_{C'}\longrightarrow C'^{(d)}
\]
be the blowups at these two diagonal points. Write $E_C,E_{C'}$ for their
exceptional divisors, $\eta_C,\eta_{C'}$ for the respective generic points,
and
\[
V_C:=\Oo_{X_C,\eta_C}\subset K(C^{(d)}),\qquad
V_{C'}:=\Oo_{X_{C'},\eta_{C'}}\subset K(C'^{(d)})
\]
for the corresponding DVRs.


\subsection{The master diagram}
\label{sec:master}

The quotient tower, compactifications, and blowups fit into the diagram
\[
\begin{tikzcd}[row sep=2.4em, column sep=2.2em]
\cP^{[d]_H}\ar[r,"H"]\ar[d]
  &(\cP/H)^{(d)}\ar[d,"q"]\ar[r,hookrightarrow,"\iota'"]
  &C'^{(d)}\ar[dd,"f^{(d)}"]
  &X_{C'}\ar[l,"\pi_{C'}"']\ar[dd,dashed,"f_X"]\\
\cP^{[d]_G}\ar[r,"H"']\ar[dr,"G"',bend right=10]
  &(\cP/H)^{[d]}\ar[d,"G/H"]
  &&\\
  &U^{(d)}\ar[r,hookrightarrow,"\iota"']
  &C^{(d)}
  &X_C\ar[l,"\pi_C"]
\end{tikzcd}
\]
The left block is the torsor tower of
\Cref{prop:torsor_tower_diagram}. The middle block records the open
immersions and compactifications. The right block records the blowups. The
rational lift $f_X$ is defined at the generic point of the exceptional
divisor and sends $\eta_{C'}$ to $\eta_C$, producing the extension of DVRs
\[
V_C\hookrightarrow V_{C'}.
\]


\subsection{The proposed induction step}

In the cyclic $p$-power setting of \Cref{sec:as_induction_motivation}, the
two intended inputs are:
\begin{description}
  \item[\textnormal{(IH$_{r-1}$)}]
    $(\cP/H)^{[d]}\to U^{(d)}$ is unramified at $\eta_C$;
  \item[\textnormal{(AS)}]
    $\cP^{[d]_H}\to(\cP/H)^{(d)}$ is unramified at $\eta_{C'}$.
\end{description}
The first is the desired contribution of induction applied to the quotient
$G/H$. The second is the desired contribution of the degree-$p$
Artin--Schreier calculation applied to the $H$-torsor over the compactified
curve $C'$. The proposed induction step asks whether these imply that
\[
\cP^{[d]_G}\longrightarrow U^{(d)}
\]
is unramified at $\eta_C$. We will see that the two assertions do not, by
themselves, control the ramification needed for this conclusion.


\subsection{A base-change lemma}

We use the following standard comparison between a Galois extension and its
base change.

\begin{lemma}\label{lem:galois-base-change}
Let $E/F$ be a finite Galois extension with group $H$, let $K/F$ be a field
extension, and form all fields inside a common overfield. Put $M:=E\cap K$.
Then restriction gives
\[
\Gal(EK/K)\cong\Gal(E/M),
\]
and, as $K$-algebras,
\[
E\otimes_FK\cong\prod_{i=1}^{[M:F]}EK.
\]
In particular, $E\otimes_FK$ is a field if and only if $E$ and $K$ are
linearly disjoint over $F$.
\end{lemma}

\begin{proof}
Write $E=F[\alpha]$ and let $g\in F[x]$ be the minimal polynomial of
$\alpha$. The translation theorem identifies $\Gal(EK/K)$ with the subgroup
of $\Gal(E/F)$ fixing $E\cap K=M$, hence with $\Gal(E/M)$.

Factor $g=g_1\cdots g_m$ into monic irreducible polynomials over $K$.
Adjoining any root of $g$ to $K$ produces $EK$, so every $g_i$ has degree
$[E:M]$. Consequently $m=[M:F]$, and
\[
E\otimes_FK\cong K[x]/(g)
  \cong\prod_{i=1}^{[M:F]}K[x]/(g_i)
  \cong\prod_{i=1}^{[M:F]}EK.
\]
The product is a field exactly when $M=F$.
\end{proof}


\section{Following the hypotheses through the valuation tower}
\label{sec:as_induction_valuation_tower}

We first make two harmless reductions. We may assume that $\cP\to U$ is
connected: replace $G$ by the image $G_0$ of its monodromy and $H$ by
$H\cap G_0$. The four quotient schemes then decompose componentwise, and the
two unramifiedness inputs pass to the connected component. We may also assume
$H\ne0$, since for $H=0$ the statement (IH$_{r-1}$) is already the desired
conclusion.

Under these reductions the four schemes in the Cartesian square of
\Cref{prop:torsor_tower_diagram} are integral, so we may pass to function
fields. Put
\[
K_C:=\Frac V_C=K(C^{(d)}),\qquad
K_{C'}:=\Frac V_{C'}=K(C'^{(d)}),
\]
and
\[
F:=K((\cP/H)^{[d]}),\qquad
E:=K(\cP^{[d]_G}).
\]
Then $F/K_C$ is Galois with group $G/H$, $E/F$ is Galois with group $H$,
and $E/K_C$ is Galois with group $G$. The Cartesian square identifies the
function field of $\cP^{[d]_H}$ with $E\cdot K_{C'}$. Its generic fibre over
$(\cP/H)^{(d)}$ is the finite domain $E\otimes_FK_{C'}$, hence a field.
By \Cref{lem:galois-base-change},
\[
E\cap K_{C'}=F,
\qquad
\Gal(EK_{C'}/K_{C'})\cong\Gal(E/F)=H.
\]

Let $V_F:=V_{C'}|_F$. Choose a place $v_E$ of $E$ above $V_F$ and a place
$v_\star$ of $EK_{C'}$ above both $v_E$ and $V_{C'}$. The relevant field
and valuation towers are
\[
\begin{tikzcd}[row sep=1.6em, column sep=1.4em]
 &EK_{C'}\arrow[dl,dash]\arrow[dr,dash,"H"]&\\
 E\arrow[dr,dash,"H"']&&K_{C'}\arrow[dl,dash]\\
 &F\arrow[d,dash,"G/H"]&\\
 &K_C&
\end{tikzcd}
\qquad
\begin{tikzcd}[row sep=1.6em, column sep=1.4em]
 &v_\star\arrow[dl,dash]\arrow[dr,dash]&\\
 v_E\arrow[dr,dash]&&V_{C'}\arrow[dl,dash]\\
 &V_F\arrow[d,dash]&\\
 &V_C&.
\end{tikzcd}
\]

We can now see exactly what each part of the proposed induction contributes:
\begin{description}
  \item[\textnormal{(IH$_{r-1}$):}]
    $F/K_C$ is unramified at $V_C$. In particular,
    $e(V_F\mid V_C)=1$.
  \item[\textnormal{(AS):}]
    $EK_{C'}/K_{C'}$ is unramified at $V_{C'}$.
  \item[Missing step:]
    one must prove that $E/F$ is unramified at $V_F$.
  \item[Goal:]
    once the missing step holds, the tower
    $K_C\subset F\subset E$ shows that $E/K_C$ is unramified at $V_C$.
\end{description}

The gap lies between the second and third items. It is measured by the
following inertia comparison.

\begin{lemma}\label{lem:inertia-injection}
Let $E/F$ be finite Galois and linearly disjoint from $K/F$. Let $w$ be a
place of $EK$ and write $w_K,w_E,w_F$ for its restrictions. Under the
restriction isomorphism $\Gal(EK/K)\cong\Gal(E/F)$ there is an injection
\[
I(w\mid w_K)\hookrightarrow I(w_E\mid w_F).
\]
After completion, the source is the inertia group of the local compositum
$\widehat E\,\widehat K/\widehat K$. In particular, if
$\widehat E\subseteq\widehat K$, then $I(w\mid w_K)=1$, regardless of the
ramification of $E/F$ at $w_F$.
\end{lemma}

\begin{proof}
Let $\sigma\in I(w\mid w_K)$. It fixes $K$, hence $F$, and restriction gives
$\sigma|_E\in\Gal(E/F)$. Because $\sigma$ preserves $w$, its restriction
preserves $w_E$. Moreover, $\sigma$ acts trivially on the residue field at
$w$, and therefore on the subfield given by the residue of $E$. Thus
$\sigma|_E\in I(w_E\mid w_F)$. Injectivity follows from the restriction
isomorphism. The completed statement is the usual identification of the
decomposition and inertia groups with those of the corresponding extension
of completed fields.
\end{proof}

Applied to the tower above, the lemma gives
\[
I(v_\star\mid V_{C'})\hookrightarrow I(v_E\mid V_F).
\]
Statement (AS) makes the group on the left trivial, whereas the missing step
requires the group on the right to be trivial. An injection with trivial
source gives no such conclusion. Surjectivity would repair the argument, but
it can fail when $K_{C'}/F$ is wildly ramified: after completion, the base
change may contain the ramified degree-$p$ layer and thereby make that layer
invisible.


\section{Failure under wild base change}
\label{sec:as_induction_wild_base_change}

The smallest possible local model already exhibits the problem. It models
the top $H\cong\Z/p\Z$ layer in the first nontrivial induction step
$G=\Z/p^2\Z$.

\begin{example}[Absorption of an Artin--Schreier layer]
\label{ex:p2-counterexample}
\label{ex:absorption-explicit}
Let $k=\bar{\F}_p$ and put $F=k((t))$. Define $K=k((\pi))$ and embed
$F\hookrightarrow K$ by
\[
t=\frac{\pi^p}{1-\pi^{p-1}}.
\]
Equivalently,
\[
t^{-1}=\pi^{-p}-\pi^{-1}
       =\wp(\pi^{-1}),
\qquad \wp(z)=z^p-z.
\]
The element $t^{-1}$ has reduced pole order one over $F$. Hence the
Artin--Schreier extension
\[
E=F(y),\qquad y^p-y=t^{-1},
\]
is totally wildly ramified of degree $p$ and has break one by
\Cref{theorem:artin_schreier_ramification}. In particular,
\[
I(E/F)=\Z/p\Z.
\]

Over $K$, however, $\pi^{-1}$ is already a solution of the same equation.
After identifying $y$ with $\pi^{-1}$ in a common overfield, one has $E=K$.
Consequently
\[
EK/K=K/K
\]
is trivial and therefore unramified. The inertia comparison is exactly
\[
\underbrace{I(EK/K)}_{1}
\hookrightarrow
\underbrace{I(E/F)}_{\Z/p\Z}.
\]

In the notation of the valuation tower, this models the possible equality of
completed fields $\widehat E=\widehat{K_{C'}}$ over $\widehat F$. The global
fields $E$ and $K_{C'}$ may still be linearly disjoint even when their chosen
completions coincide. Thus the example disproves the general local inference
``unramified after base change implies unramified before base change.'' It
does not assert that this particular extension $K/F$ is produced by the map
$q$ for a specified global symmetric-power torsor.
\end{example}

The mechanism is visible in the valuations. Over $F$, the pole order of
$t^{-1}$ is one and cannot be reduced by an element of $F$. Under the wildly
ramified extension $K/F$, that pole order becomes $p$, and $K$ supplies the
element $\pi^{-1}$ of intermediate valuation needed for the
Artin--Schreier change of variable. The base change does not show that the
original ramification was absent; it absorbs it.


\section{What the example proves}
\label{sec:as_induction_takeaway}

The example does not contradict part~(3) of
\Cref{theorem:cyclic_cases_after_blowup}. That theorem is proved in Chapter~2
by a direct Artin--Schreier--Witt calculation. Rather, the example isolates
the missing implication in the proposed alternative proof:
\begin{enumerate}
  \item the induction hypothesis for $G/H$ controls the lower extension
        $F/K_C$;
  \item the degree-$p$ Artin--Schreier calculation controls the top layer
        only after base change to $K_{C'}$;
  \item a wild base change can absorb the ramification of $E/F$, so
        unramifiedness of $EK_{C'}/K_{C'}$ does not recover unramifiedness of
        $E/F$;
  \item repairing the induction would require additional geometric control
        of the comparison extension $K_{C'}/F$, strong enough to prevent
        this absorption;
  \item the direct Witt-vector calculation avoids this descent problem by
        computing the full $\Z/p^r\Z$-torsor at the exceptional valuation in
        one step.
\end{enumerate}

Thus the obstruction is not the decomposition of the torsor itself. The
obstruction is that the two induction inputs measure the top layer on opposite
sides of an uncontrolled wild base change.
